% Part III: CLIF. Sources: medorama/src/formats/clif.rs (its module doc and code are the
% authority: tables(), rank(), import_folder(), the converters), docs/administration.md (the
% clif- row of "Import formats"); main.rs (import-clif), server.rs (tables/clif),
% web/index.html (the CLIF folder box).
\chapter{CLIF}
\label{chap:clif}

This chapter explains how Medorama imports CLIF databases. It is written for data managers who
hold intensive-care data in CLIF. CLIF, the \emph{Common Longitudinal ICU data Format}, is an
open format for the data of intensive care: a set of tables (patients, hospitalizations, ADT,
vital signs, laboratory results, respiratory support, medications, assessments, microbiology and
others) with fixed columns and standard categories, so that the data of different hospitals read
alike. Its data dictionary and tools are published by the CLIF consortium under the Apache-2.0
license: \url{https://github.com/Common-Longitudinal-ICU-data-Format/CLIF}. Medorama reads every
table of CLIF 2.1, and five tables of CLIF 3.0 (\code{input}, \code{output},
\code{intermittent\_dialysis}, \code{validated\_diagnosis} and \code{clinical\_notes\_facts},
with their columns as in clifpy's 3.0 schemas).

Medorama reads CLIF as tables. Each CLIF table, as CSV or Parquet, becomes a database of the
catalog, as an imported table does (\cref{chap:importing}), through one import format per table.
CLIF holds no waveforms, and it is not converted to WFDB.

\section{How CLIF is read}
\label{sec:clif-reading}

\begin{itemize}
\item \textbf{Times.} CLIF's times are in UTC, written with their offset
  (\code{YYYY-MM-DD HH:MM:SS+00:00}), and they are read as such, whatever the catalog's time
  zone. A CLIF table's time written without an offset is read as UTC too, not in the catalog's
  time zone.
\item \textbf{Names.} CLIF's standard categories (its mCIDE: \code{vital\_category},
  \code{lab\_category}, \code{med\_category}, \code{location\_category} and so on) name the
  entries, written as words (\code{heart\_rate} becomes \code{Heart rate}), so that every site's
  data read alike. The sites' raw names (\code{vital\_name}, \code{lab\_name}, \code{med\_name}
  and so on) stay, as the database's own columns.
\item \textbf{Units.} Each entry has its unit: CLIF's fixed units for vital signs, settings and
  rates (\cref{sec:clif-units}), the reference unit of a laboratory result, or the dose unit of a
  drug.
\item \textbf{Patients and stays.} \code{patient\_id} names the patient and
  \code{hospitalization\_id} the encounter. Most tables name only the hospitalization, and their
  entries go to its patient.
\item \textbf{Rows left out} are counted by reason, and never named, because rows hold
  identifiers. A row without the keys that its table needs (\code{rows without
  hospitalization\_id}), or one that holds no value to keep, makes no entry.
\end{itemize}

\section{Importing a CLIF folder}
\label{sec:clif-importing}

A CLIF database is a folder of tables, with one file per table (or several files per table).
Medorama imports such a folder in one step:
\begin{enumerate}
\item Every file directly in the folder whose name ends in \file{.csv} or \file{.parquet} is read
  for its columns. Subfolders are not read.
\item Each file is recognized as a CLIF table by its columns (\cref{sec:clif-detect}). A file
  that is not a CLIF table that Medorama reads is named in the report and left out.
\item Each table becomes a database of its own, named after it (\code{clif\_patient},
  \code{clif\_vitals}, \code{clif\_medication\_admin} and so on), with the source label
  \code{clif}. A table in several files (a year per file, or CLIF's continuous and intermittent
  medications) is one database.
\item The tables are imported in order (\cref{sec:clif-order}), each read row by row.
\end{enumerate}

\paragraph{In the page.} In \ui{Import} $\rightarrow$ \ui{Databases}, in the card \ui{Add a
database}, choose \ui{Import tables}, open \ui{A CLIF folder on the server}, and type the
folder's path on the server (\file{/data/clif}), or click \ui{Choose\ldots} and open it. Then click \ui{Preview}, a dry run that reads
and counts everything and keeps nothing, or \ui{Import the folder}. The folder must be inside
\code{data\_roots}. The import runs in the background, with its progress shown like any other.
The report gives, for each table, its database's name, its files and what it added, and then
the files that are not CLIF tables.

\paragraph{On the command line,} \cmd{medorama import-clif} imports a folder:
\begin{lstlisting}[style=signature]
medorama import-clif DIR [--source LABEL] [--dry-run]
\end{lstlisting}
\begin{params}
\item[DIR] The CLIF folder.
\item[--source LABEL] The source label of every table's rows. Default: \code{clif}.
\item[--dry-run] Check and count only; store nothing.
\end{params}
The following commands preview the import of a folder and then import it:
\begin{lstlisting}[style=shell]
medorama import-clif /data/clif --dry-run
medorama import-clif /data/clif
\end{lstlisting}
The command prints a line per table: its database, its files, its rows, the patients,
encounters, bed stays and items it added, and the rows held back, with up to five of its
warnings. Unlike the page and the API, the command does not check the folder against
\code{data\_roots}, because it runs as the operating-system user who can read the home folder
(\cref{chap:cli}).

\paragraph{Through the API,} an administrator sends \code{POST api/v1/tables/clif} with the JSON
body \code{\{"dir": "/data/clif", "source": "clif", "dry\_run": false, "background": true\}}. The
folder must be inside \code{data\_roots}. With \code{background}, the import runs as a task that
the page follows (\cref{chap:http-api}).

\paragraph{One table at a time.} A single CLIF table can also be imported with
\cmd{medorama import-table}. It is recognized by its columns, as any known export is, or named
with \opt{-{}-format clif-vitals} (\cref{chap:importing}). Import the patients, then the
hospitalizations, then the ADT, before the rest. Susceptibilities cannot be imported alone
(\cref{sec:clif-micro}).

\section{How tables are told apart}
\label{sec:clif-detect}

A file is a CLIF table when it has all the columns of \cref{tab:clif-detect} for that table, and
none of the columns that belong to another table that shares them. Column names are compared
ignoring case and surrounding spaces. The formats are those whose key begins with \code{clif-};
\cmd{medorama formats} lists them.

\begin{xltabular}{\linewidth}{@{}P{4.2cm}P{5.2cm}L@{}}
\caption{The columns that recognize each CLIF table.}
\label{tab:clif-detect}\\
\toprule
\tabhead Format & \tabhead Columns it needs & \tabhead Unless it has \\
\midrule
\endfirsthead
\toprule
\tabhead Format & \tabhead Columns it needs & \tabhead Unless it has \\
\midrule
\endhead
\bottomrule
\endfoot
\code{clif-patient} & \code{patient\_id}, \code{sex\_category}, \code{birth\_date} &
  \code{hospitalization\_id} \\
\code{clif-hospitalization} & \code{patient\_id}, \code{hospitalization\_id},
  \code{admission\_dttm} & \\
\code{clif-adt} & \code{hospitalization\_id}, \code{in\_dttm}, \code{location\_category} & \\
\code{clif-vitals} & \code{hospitalization\_id}, \code{vital\_category}, \code{vital\_value} & \\
\code{clif-invasive-hemodynamics} & \code{hospitalization\_id}, \code{measure\_category},
  \code{measure\_value} & \\
\code{clif-intake-output} & \code{hospitalization\_id}, \code{in\_out\_flag}, \code{amount} & \\
\code{clif-input} & \code{hospitalization\_id}, \code{input\_category}, \code{input\_volume} & \\
\code{clif-output} & \code{hospitalization\_id}, \code{output\_category}, \code{output\_value} &
  \\
\code{clif-labs} & \code{hospitalization\_id}, \code{lab\_category} & \\
\code{clif-respiratory-support} & \code{hospitalization\_id}, \code{device\_category},
  \code{recorded\_dttm} & \code{mcs\_group} \\
\code{clif-ecmo-mcs} & \code{hospitalization\_id}, \code{mcs\_group} & \\
\code{clif-crrt-therapy} & \code{hospitalization\_id}, \code{crrt\_mode\_category} & \\
\code{clif-intermittent-dialysis} & \code{hospitalization\_id}, \code{net\_ultrafiltration\_out}
  & \\
\code{clif-medication-admin} & \code{hospitalization\_id}, \code{med\_category},
  \code{admin\_dttm} & \\
\code{clif-medication-orders} & \code{hospitalization\_id}, \code{med\_order\_id},
  \code{order\_start\_dttm} & \\
\code{clif-transfusion} & \code{hospitalization\_id}, \code{component\_name},
  \code{transfusion\_start\_dttm} & \\
\code{clif-patient-assessments} & \code{hospitalization\_id}, \code{assessment\_category} & \\
\code{clif-position} & \code{hospitalization\_id}, \code{position\_category} & \\
\code{clif-code-status} & \code{patient\_id}, \code{code\_status\_category} & \\
\code{clif-key-icu-orders} & \code{hospitalization\_id}, \code{order\_category},
  \code{order\_dttm} & \\
\code{clif-therapy-details} & \code{hospitalization\_id}, \code{therapy\_element\_category} & \\
\code{clif-provider} & \code{hospitalization\_id}, \code{provider\_id} & \\
\code{clif-place-based-index} & \code{hospitalization\_id}, \code{index\_name},
  \code{index\_value} & \\
\code{clif-clinical-trial} & \code{hospitalization\_id}, \code{trial\_id} & \\
\code{clif-microbiology-culture} & \code{organism\_category}, \code{collect\_dttm} &
  \code{result\_category}, \code{micro\_order\_name} \\
\code{clif-microbiology-nonculture} & \code{organism\_category}, \code{result\_category} & \\
\code{clif-microbiology-susceptibility} & \code{organism\_id}, \code{antimicrobial\_category} &
  \\
\code{clif-patient-procedures} & \code{hospitalization\_id}, \code{procedure\_code} & \\
\code{clif-hospital-diagnosis} & \code{hospitalization\_id}, \code{diagnosis\_code} &
  \code{start\_dttm}, \code{source\_type} \\
\code{clif-patient-diagnosis} & \code{patient\_id}, \code{diagnosis\_code}, \code{start\_dttm} &
  \\
\code{clif-validated-diagnosis} & \code{hospitalization\_id}, \code{diagnosis\_category},
  \code{validation\_method} & \\
\code{clif-clinical-notes-facts} & \code{hospitalization\_id}, \code{note\_id},
  \code{note\_type\_category} & \\
\end{xltabular}

\section{The import order}
\label{sec:clif-order}

The tables of a folder are imported in this order, and in the order of their formats' keys within
each step:
\begin{enumerate}
\item \code{patient}, which says who the patients are;
\item \code{hospitalization}, which says which stays there are, and of which patient;
\item \code{adt}, the bed stays, which name only their hospitalization;
\item every other table, whose rows name their hospitalization (or patient);
\item \code{microbiology\_susceptibility}, last, after the cultures it names.
\end{enumerate}
In this way, every entry finds its encounter at once. A table imported out of this order is not
lost: its items wait for their patient or encounter, and are matched when the table that names
it comes (\cref{chap:importing}). ADT rows are the exception: a bed stay that names only its
hospitalization needs the hospitalization first.

\section{What each table becomes}
\label{sec:clif-tables}

\Cref{tab:clif-tables} lists every CLIF table Medorama reads, its category
(\cref{chap:importing}) and what its rows become: the entry's name, its value and unit, and the
columns kept as the database's own.

\begin{xltabular}{\linewidth}{@{}P{3.6cm}P{2.0cm}L@{}}
\caption{The CLIF tables Medorama reads, and what each becomes.}
\label{tab:clif-tables}\\
\toprule
\tabhead CLIF table & \tabhead Category & \tabhead Becomes \\
\midrule
\endfirsthead
\toprule
\tabhead CLIF table & \tabhead Category & \tabhead Becomes \\
\midrule
\endhead
\bottomrule
\endfoot
\code{patient} & (patients) & The patient, by \code{patient\_id}: the sex from
  \code{sex\_category} (\code{Male} \code{M}, \code{Female} \code{F}), the date of birth (its first
  ten characters); race, ethnicity, language and the time of death as columns \\
\code{hospitalization} & (encounters) & An encounter, by \code{hospitalization\_id}: admission
  (\code{admission\_dttm}), discharge, disposition (\code{discharge\_category}); the admission
  type as a column \\
\code{adt} & (bed stays) & A bed stay of the hospitalization, its unit the
  \code{location\_category} (\code{icu} shows as ICU), from \code{in\_dttm} to \code{out\_dttm};
  the hospital and the location's name as columns. No bed, so ADT does not link recordings by bed
  and time \\
\code{vitals} & Vitals & Each vital sign, by \code{vital\_category}, in CLIF's units
  (\cref{tab:clif-units}); rows without a number are left out \\
\code{invasive\_hemodynamics} & Vitals & Each measurement, by \code{measure\_category}: CVP, right
  atrial and ventricular pressures, PA systolic, diastolic and mean pressures, PCWP, cardiac
  output by thermodilution and by Fick \\
\code{intake\_output} (2.1) & Vitals & Each intake or output, \code{Intake: <fluid>} or
  \code{Output: <fluid>} (by \code{in\_out\_flag} 1 or 0, the fluid's name), in mL \\
\code{input}, \code{output} (3.0) & Vitals & \code{Intake: <category>}, in mL;
  \code{Output: <category>} (urine, drains \dots), in \code{output\_unit} (mL, or occurrences:
  emesis \dots); the group and the raw name as columns \\
\code{labs} & Labs & Each result, by \code{lab\_category}: the numeric value
  (\code{lab\_value\_numeric}), else the text (\code{lab\_value}), in the
  \code{reference\_unit}; the LOINC code; at its collection time, else its result's, else its
  order's; the raw name, specimen and order and result times as columns \\
\code{respiratory\_support} & Life support & An entry \code{Respiratory support} (the device),
  and each setting and observation of the row an entry of its own (\cref{tab:clif-units}); the
  mode and tracheostomy as columns \\
\code{ecmo\_mcs} & Life support & \code{Circulatory support} (the device, or the MCS group), blood
  flow, sweep gas flow and FdO2, and the device's work rate (RPM, Impella power \dots) named by its
  control parameter's category, else \code{Device rate} \\
\code{crrt\_therapy} & Life support & \code{CRRT mode} (\code{CVVHDF} \dots) and each rate
  \\
\code{intermittent\_dialysis} (3.0) & Life support & Each rate, named \code{\dots{} (IHD)} \\
\code{medication\_admin\_continuous}, \code{medication\_admin\_intermittent} & Medications & Each
  administration, by \code{med\_category}, with its dose and dose unit; a dose not given is the
  event without a dose (\cref{sec:clif-not-given}); the group, route, action and raw name as
  columns \\
\code{medication\_orders} & Medications & \code{<drug> order}: dose and unit, from start to end;
  route, status, frequency, PRN and order time as columns \\
\code{transfusion} & Transfusions & The component (\code{component\_name}), its volume and unit,
  its ISBT 128 product code \\
\code{patient\_assessments} & Assessments & Each assessment (GCS, RASS, CAM \dots), by
  \code{assessment\_category}: its numerical, else categorical, else text value \\
\code{position} & Assessments & \code{Position}: prone or not prone \\
\code{code\_status} & Assessments & \code{Code status}, of the patient, from its start \\
\code{key\_icu\_orders} & Assessments & Each order (physical and occupational therapy \dots), by
  \code{order\_category}, with its status \\
\code{therapy\_details} & Assessments & Each element of a therapy session and its value \\
\code{provider} & Assessments & \code{Care team}: each role, from start to stop; never who
  (\code{provider\_id} is not read) \\
\code{place\_based\_index} & Assessments & Each index (Area Deprivation Index, Social
  Vulnerability Index \dots) of the hospitalization, with its version \\
\code{clinical\_trial} & Assessments & \code{Clinical trial}: the trial, at randomization (else
  consent); the arm, consent, randomization and withdrawal as columns \\
\code{microbiology\_culture} & Microbiology & A culture, named by its method and fluid
  (\code{Culture blood}), its value the organism; the LOINC code; its \code{organism\_id} kept
  (\cref{sec:clif-micro}) \\
\code{microbiology\_nonculture} & Microbiology & Each organism tested (PCR \dots), detected or
  not, with its unit and range \\
\code{microbiology\_susceptibility} & Microbiology & \code{<organism>: <antimicrobial>}, its
  result (susceptible, resistant \dots), on its culture (\cref{sec:clif-micro}) \\
\code{patient\_procedures} & Procedures & Each billed procedure, named by its code, in its
  vocabulary (\code{procedure\_code\_format}) \\
\code{hospital\_diagnosis} & Conditions & Each diagnosis of the hospitalization, by its code and
  vocabulary, without a time; primary and present on admission as columns \\
\code{patient\_diagnosis} & Conditions & Each diagnosis from clinical documentation (problem
  list, history), by its code, from start to end \\
\code{validated\_diagnosis} (3.0) & Conditions & Each clinician-validated label and its status
  (confirmed, ruled out, uncertain); never who reviewed it \\
\code{clinical\_notes\_facts} (3.0) & Notes & Each note's revision, by note type: its time, type,
  author's specialty, status and revision, without text \\
\end{xltabular}

\subsection{Names and units}
\label{sec:clif-units}

\begin{xltabular}{\linewidth}{@{}P{5.0cm}P{5.2cm}L@{}}
\caption{The names and units of CLIF's vital signs, measurements, settings and rates.}
\label{tab:clif-units}\\
\toprule
\tabhead CLIF category or column & \tabhead Name & \tabhead Unit \\
\midrule
\endfirsthead
\toprule
\tabhead CLIF category or column & \tabhead Name & \tabhead Unit \\
\midrule
\endhead
\bottomrule
\endfoot
\multicolumn{3}{@{}l}{\textbf{vitals: \code{vital\_category}}} \\
\code{temp\_c} & Temperature & \textdegree C \\
\code{heart\_rate} & Heart rate & bpm \\
\code{sbp}, \code{dbp}, \code{map} & Systolic BP, Diastolic BP, Mean arterial pressure & mmHg \\
\code{spo2} & SpO2 & \% \\
\code{respiratory\_rate} & Respiratory rate & /min \\
\code{height\_cm}, \code{weight\_kg} & Height, Weight & cm, kg \\
\addlinespace
\multicolumn{3}{@{}l}{\textbf{invasive\_hemodynamics: \code{measure\_category}}} \\
\code{cvp}, \code{ra}, \code{rv} & CVP, Right atrial pressure, Right ventricular pressure & mmHg \\
\code{pa\_systolic}, \code{pa\_diastolic}, \code{pa\_mean} & PA systolic, PA diastolic, PA mean
  pressure & mmHg \\
\code{pcwp} & PCWP & mmHg \\
\code{cardiac\_output\_thermodilution}, \code{cardiac\_output\_fick} & Cardiac output
  (thermodilution), Cardiac output (Fick) & L/min \\
\addlinespace
\multicolumn{3}{@{}l}{\textbf{respiratory\_support: columns}} \\
\code{fio2\_set} & FiO2 (set) & fraction \\
\code{lpm\_set}, \code{flow\_rate\_set} & O2 flow (set), Flow rate (set) & L/min \\
\code{tidal\_volume\_set}, \code{tidal\_volume\_obs} & Tidal volume (set), Tidal volume & mL \\
\code{resp\_rate\_set}, \code{resp\_rate\_obs} & Respiratory rate (set), Respiratory rate
  (ventilator) & /min \\
\code{pressure\_control\_set}, \code{pressure\_support\_set} & Pressure control (set), Pressure
  support (set) & cmH2O \\
\code{peak\_inspiratory\_pressure\_set}, \code{\dots\_obs} & Peak inspiratory pressure (set),
  Peak inspiratory pressure & cmH2O \\
\code{peep\_set}, \code{peep\_obs} & PEEP (set), PEEP & cmH2O \\
\code{plateau\_pressure\_obs}, \code{mean\_airway\_pressure\_obs} & Plateau pressure, Mean airway
  pressure & cmH2O \\
\code{inspiratory\_time\_set} & Inspiratory time (set) & s \\
\code{minute\_vent\_obs} & Minute ventilation & L/min \\
\addlinespace
\multicolumn{3}{@{}l}{\textbf{ecmo\_mcs: columns}} \\
\code{flow} & Blood flow (MCS) & L/min \\
\code{sweep\_set} (or \code{sweep}) & Sweep gas flow (set) & L/min \\
\code{fdo2\_set} (or \code{fdo2}) & FdO2 (set) & fraction \\
\addlinespace
\multicolumn{3}{@{}l}{\textbf{crrt\_therapy: columns}} \\
\code{blood\_flow\_rate} & Blood flow rate & mL/min \\
\code{pre\_filter\_replacement\_fluid\_rate}, \code{post\_filter\_\dots} & Pre-filter, Post-filter
  replacement fluid rate & mL/h \\
\code{dialysate\_flow\_rate} & Dialysate flow rate & mL/h \\
\code{ultrafiltration\_out} & Ultrafiltration & mL/h \\
\addlinespace
\multicolumn{3}{@{}l}{\textbf{intermittent\_dialysis: columns}} \\
\code{blood\_flow\_rate} & Blood flow rate (IHD) & mL/min \\
\code{dialysate\_flow\_rate}, \code{net\_ultrafiltration\_out} & Dialysate flow rate (IHD), Net
  ultrafiltration (IHD) & none (clifpy's 3.0 schema gives none) \\
\end{xltabular}

A category not in these lists is named by its words (\code{etco2} becomes \code{Etco2}), without a
unit.

\section{Susceptibilities on their cultures}
\label{sec:clif-micro}

A row of \code{microbiology\_susceptibility} names only its organism, by \code{organism\_id}, and
not its patient, its hospitalization or its time. When a folder is imported, Medorama therefore
first reads the folder's \code{microbiology\_culture} files and keeps each culture by its
\code{organism\_id}. It then puts each susceptibility on its culture, with the culture's patient,
encounter and collection time. The susceptibility is named \code{<organism>: <antimicrobial>}
(\code{Escherichia coli: Ceftriaxone}), with the culture and the organism as columns. A
susceptibility whose culture is not in the folder, and every row of a susceptibility table
imported on its own, is left out and counted (\code{rows without their culture}).

\section{Doses not given}
\label{sec:clif-not-given}

An administration that was not given is kept as an event, without a dose and unit, so that it is
never counted or searched as a dose given. CLIF marks such an administration with the
\code{mar\_action\_group} \code{not\_administered} (stopped, paused, held) or, without a group,
with the \code{mar\_action\_category} \code{stop}. Its action stays in the \code{action} column.

\section{What is left out, and why}
\label{sec:clif-left-out}

A format keeps only the columns it names, and everything else of a CLIF table is left out. The
following are left out on purpose:
\begin{itemize}
\item \textbf{The geography} of \code{hospitalization} (ZIP codes, census blocks), because these
  are identifiers under HIPAA.
\item \textbf{Devices' serial IDs} (\code{device\_id}), for the same reason.
\item \textbf{\code{age\_at\_admission}}, because ages are computed from the date of birth, and
  capped in Privacy Mode.
\item \textbf{Staff IDs} (\code{provider\_id}, \code{author\_id}, \code{cosigner\_id},
  \code{reviewer\_id}), because formats keep no staff.
\item \textbf{\code{lda\_id}} (lines, drains and airways), which is not read.
\item \textbf{Note text}, which CLIF does not hold. \code{clinical\_notes\_facts} gives a note's
  facts (its type, time, status and revision), which become entries of the Notes column without
  text.
\end{itemize}

\begin{phibox}
A CLIF database still holds identifiers: \code{patient\_id}, \code{hospitalization\_id}, dates of
birth and every time. Once imported, they are kept like any other table's identifiers
(\cref{chap:importing}). They are shown only to users with the PHI permission out of Privacy
Mode, and are otherwise replaced by pseudonyms and moved dates (\cref{chap:protecting}).
\end{phibox}
